\section{Discussion}
% wordcount 510
% Restate the Research Question
In this study, we proposed a novel perspective on how perceptual information is embedded in correlated neural variability and how this covariance gets involved in neural computation.
We introduced an encoding scheme using the covariance of sensory neurons to capture complex perceptual wholes, a significant advance considering neurons' stochastic activities in the time domain.
Through a motion direction detection task, we showed that downstream networks are capable of learning covariance-based computation to decode this encoded information effectively. 
This process allows for the transformation of perceptual information from the covariance observed in upstream neurons to the firing rates of downstream neurons, based on the perceptual disentangling hypothesis~\cite{DiCarlo2007, DiCarlo2012}. 
Additionally, we demonstrated that SNNs could perform this computation implicitly, leading to lossless recovery of direction information. 
Moreover, our approach proved advantageous in a complex natural image classification task, enhancing model performance and inference speed, underscoring the significance of correlated neural variability in sophisticated information processing.

% Theoretical and Practical Implications
Our encoding scheme hinges on two essential components: a variety of neurons responsive to different stimulus aspects and an appropriate observation window to track stimulus-induced changes in neuronal responses. 
This approach aligns with the concept of a combinatorial code~\cite{Malnic1999, Osborne2008}, where neural responses at different timescales represent different stimulus characteristics~\cite{Panzeri2010}.
The resulting activity patterns of neural populations, shaped by stimuli, form a consistent neural manifold~\cite{DiCarlo2007, DiCarlo2012}, which remains stable despite the variability of the trial-to-trial response, as encapsulated by the covariance in our scheme.
From the decoding perspective, downstream neurons serve dual roles as integrators and coincidence detectors~\cite{Koenig1996}, because processing perceptual information requires extended timescales and collaboration between sensory neurons. 
We describe this process as covariance-based computation, where the covariance of sensory neuron spike trains dominates the firing rates of downstream neurons.
This finding suggests a hierarchical structure in neural processing, achieved through iterative covariance-based computations. 
As one ascends this hierarchy, the timescale and complexity of the encoded information increase, leading to a progressively linear representation of perceptual information~\cite{piasini2021temporal, Siegle2021}.

%Limitations and Future Research
While our results are promising, we must acknowledge certain limitations. 
The downstream networks were trained using gradient-based backpropagation, capable of explicitly computing gradients related to covariance. 
However, the exact learning mechanisms of the brain, especially akin to backpropagation, are still not fully understood.  
Further research could delve deeper into the effects of local learning rules that are impacted by correlated neural activities.
For example, employing a burst-dependent synaptic plasticity rule~\cite{Payeur2021} could modify error signals in feedback connections. 
Such exploration could shed light on the prevalence of noise correlation in the brain and its potential role in enhancing learning processes.

%Conclusion
The approach developed in this work also has broader implications for SNN training. 
We have demonstrated the benefits of incorporating second-order information, such as covariance, into model training. 
Integrating covariance into ANNs requires a quadratic number of parameters, making it difficult to use for larger models~\cite{Wang2020}. 
In comparison, SNNs can unfold the covariance over time, which allows them to implicitly process the correlation between variables, thus improving performance without the need for extra parameters.
Ultimately, our work paves the way for novel studies on the roles of noise and neural correlation in the brain, as well as the development of powerful learning algorithms for SNNs.
